\documentclass[spanish,letterpaper,oneside]{article}
\def\documenttitle {Fundamento y exigibilidad de los derechos humanos}
\def\documentsubtitle {Participacion sobre teorias de fundamentacion}
\def\documentsubject {Actividad 4 - Antecedentes de los derechos humanos}
\def\documentauthor {Martin Jonathan de la Cruz}
\def\coursename {Antecedentes de los derechos humanos}
\def\coursecode {005}
\def\universityname {Universidad Abierta y a Distancia de Mexico}
\def\universityfaculty {Licenciatura en Derecho}
\def\universitydepartment {Antecedentes de los derechos humanos}
\def\universitydepartmentimage {departamentos/UnADM}
\def\universitydepartmentimagecfg {height=1.57cm}
\def\universitylocation {Roma Norte, Ciudad de Mexico}
\def\authortable {\begin{tabular}{ll}
Alumno: & Martin Jonathan de la Cruz \\
Matricula: & ES2611202040 \\
Actividad interna: & 4 \\
Semana: & 2 \\
Figura docente: & Daniel Alexis Lozano Ortega \\
Estado: & Propuesta no publicada \\
Fecha: & \today
\end{tabular}}
\input{template}
\usepackage{helvet}
\renewcommand{\familydefault}{\sfdefault}
\usepackage[most]{tcolorbox}
\newtcolorbox{forobox}{enhanced,breakable,colback=white,
colframe=green!30!black,leftrule=2pt,boxrule=0.6pt,arc=2pt}
\begin{document}
\templatePortrait
\templatePagecfg
\templateFinalcfg
\onehalfspacing
\section{Introduccion}
La pregunta del foro exige elegir una fundamentacion y explicar su
importancia actual. Unidad 1 distingue derechos inherentes a la persona,
reconocimiento normativo y mecanismos de proteccion. El texto siguiente
conserva la propuesta local para revision personal; no acredita publicacion.
\section{Dignidad y garantias institucionales}
La consigna oficial, Actividad 2 de semana 2, pide una unica
participacion principal de 50 a 100 palabras; no exige replicas.
Su destino es el modulo Moodle 211997.
\begin{forobox}
El iusnaturalismo ofrece una explicacion convincente del origen de los derechos humanos porque los vincula con la dignidad de cada persona, no con una concesion del gobernante. Ese fundamento permite cuestionar normas injustas y defender a quienes han sido excluidos. Sin embargo, la justificacion moral necesita reconocimiento juridico y garantias efectivas: sin instituciones y mecanismos de exigibilidad, la dignidad puede quedarse en una proclamacion. Por ello, considero importante distinguir el fundamento del derecho de las condiciones que permiten protegerlo en la sociedad actual.
\end{forobox}
La propuesta contiene 83 palabras y debe confirmarse como postura del
estudiante. El apoyo conceptual procede del cuadernillo de Unidad 1,
recurso 211995; este parrafo documental no forma parte del texto que
se publicaria. No se anaden citas literales ni requisitos genericos de
foro que contradigan el limite y la consigna especificos.
\section{Conclusiones}
La propuesta vincula dignidad y exigibilidad sin confundir sus funciones.
Antes de publicar, corresponde revisar su fidelidad a Unidad 1 y
comprobar si ya existe una participacion propia para evitar duplicarla.
Este reporte documenta preparacion, no envio.\footnote{GitHub Copilot
apoyo la redaccion y organizacion; la postura requiere revision personal.}
\section*{Fuente y destino}
Lozano Ortega, D. A. (2026). \emph{Unidad 1. Bases teoricas basicas
relacionadas con los derechos humanos en Mexico}. UnADM.
\url{https://aulavirtual.unadmexico.mx/mod/resource/view.php?id=211995}.

Foro: \url{https://aulavirtual.unadmexico.mx/mod/forum/view.php?id=211997}.
\end{document}